\subsection{A learning-free recipe and what its budget buys}
\label{sec:recipe}

If learning does not pay, what does? The evidence points to two cheap ingredients: a strong start and a destruction size chosen \emph{offline}.

\paragraph{A strong start}
LDM, the multiway differencing heuristic \citep{KarmarkarKarp1983,MichielsEtal2007}, improves on LPT (mean GAP against the certified bound $1.23\%$ against $2.18\%$) but is worse than MULTIFIT ($0.43\%$; MULTIFIT wins $585$ of the
$725$ decisive synthetic instances). Plain \texttt{ig\_ql}, which starts from LPT, loses to LDM in $21\%$ of the synthetic and $36\%$ of the published instances. Starting the IG from the best of LPT, MULTIFIT and LDM
(\textsc{IG-QL}$^{+}$) recovers the large-$n$ classes: it ties the best constructive heuristic in every instance with $n\geq100$, improves on it in $19\%$ (synthetic) and $34\%$ (published) of the instances and is never worse
(Table~\ref{tab:f12_stronger}; mean GAP $0.17\%$ against $0.35\%$ for \texttt{ig\_ql} and $0.43\%$ for MULTIFIT; $69.3\%$ provably optimal against $34.3\%$ and $48.5\%$). Even with the strong start the policy adds
almost nothing: $-0.010\%$ (synthetic) and $-0.013\%$ (published) against random destruction from the same seed, inside the $\pm0.05\%$ margin. \textsc{IG-QL}$^{+}$ lies on the quality--time Pareto frontier in all three size ranges
(Figure~\ref{fig:pareto}), at $0.2$--$0.7$ seconds per run; MULTIFIT costs $50$--$500$ times less than \texttt{ig\_ql}.

% Auto-generated by analysis/make_table_stronger_baselines.py (Entrada 102)
\begin{table}[t]
\centering
\caption{Stronger baselines against the certified arc-flow lower bound (LB$^{\mathrm{AF}}$): mean GAP (\%) and share of instances solved \emph{provably} to optimality (makespan $=$ LB$^{\mathrm{AF}}$). \textsc{IG-QL}$^{+}$ is \textsc{IG-QL} started from the best of LPT, \textsc{MultiFit} and \textsc{LDM}. Classes sorted by $n$; $N=100$ per class ($30$ for $(20,4)$). Bold: lowest GAP in the row (ties all bold).}
\label{tab:f12_stronger}
{\scriptsize\setlength{\tabcolsep}{3.5pt}
\begin{tabular}{@{}rr rrrr rrrr@{}}
\toprule
 & & \multicolumn{4}{c}{GAP vs LB$^{\mathrm{AF}}$ (\%)} & \multicolumn{4}{c}{Provably optimal (\%)} \\
\cmidrule(lr){3-6}\cmidrule(lr){7-10}
$n$ & $m$ & \textsc{MultiFit} & \textsc{LDM} & \textsc{IG-QL} & \textsc{IG-QL}$^{+}$ & \textsc{MultiFit} & \textsc{LDM} & \textsc{IG-QL} & \textsc{IG-QL}$^{+}$ \\
\midrule
20 & 4 & 0.759 & 0.493 & \textbf{0.000} & 0.026 & 10 & 27 & 100 & 93 \\
20 & 6 & 1.454 & 2.231 & \textbf{0.169} & 0.173 & 10 & 1 & 74 & 73 \\
30 & 8 & 1.123 & 1.789 & 0.211 & \textbf{0.165} & 6 & 3 & 59 & 68 \\
40 & 5 & 0.199 & 0.044 & 0.021 & \textbf{0.009} & 45 & 82 & 91 & 96 \\
50 & 10 & 0.390 & 0.692 & 0.291 & \textbf{0.243} & 24 & 12 & 31 & 42 \\
60 & 20 & 0.851 & 5.550 & 0.820 & \textbf{0.775} & 9 & 2 & 12 & 9 \\
75 & 15 & 0.239 & 0.898 & 0.559 & \textbf{0.223} & 44 & 4 & 4 & 46 \\
100 & 20 & \textbf{0.175} & 1.047 & 0.666 & \textbf{0.175} & 56 & 1 & 1 & 56 \\
150 & 15 & 0.039 & 0.026 & 0.166 & \textbf{0.016} & 80 & 87 & 21 & 92 \\
150 & 30 & \textbf{0.101} & 1.480 & 0.816 & \textbf{0.101} & 74 & 0 & 0 & 74 \\
250 & 25 & \textbf{0.002} & 0.026 & 0.175 & \textbf{0.002} & 99 & 87 & 21 & 99 \\
300 & 20 & 0.003 & 0.004 & 0.074 & \textbf{0.000} & 98 & 97 & 44 & 100 \\
\midrule
\multicolumn{2}{c}{All (1{,}130)} & 0.425 & 1.233 & 0.351 & \textbf{0.167} & 48.5 & 34.0 & 34.3 & 69.3 \\
\bottomrule
\end{tabular}}
\end{table}


\begin{figure}[t]
\centering
\includegraphics[width=0.92\textwidth]{fig_pareto.pdf}
\caption{Quality--time trade-off on the synthetic classes, in three size ranges (each marker is the mean over the four classes of the panel; GAP against the certified bound; CPU time is the geometric mean of the
class medians, log scale, from one sequential session on an idle machine). The dashed line is the Pareto frontier. \textsc{IG-QL}$^{+}$: \texttt{ig\_ql} started from the best of LPT, MULTIFIT and LDM.}
\label{fig:pareto}
\end{figure}

\paragraph{A destruction size chosen offline}
\textsc{IG-Tuned}$^{+}$ combines both ingredients: Iterated Greedy with
random destruction, $d/n$ chosen per class on $30$ separate instances (per family on $553$ published instances never run, for the external data), started from the best of LPT, MULTIFIT and LDM,
with the budget $T=1{,}000$ and no learned component. Table~\ref{tab:f22_igcv} gives the pre-specified comparison (all eight tests confirmed). The recipe beats MULTIFIT by
$0.29\%$ (synthetic) and $0.85\%$ (published) and LDM by about $1.1\%$, never returns a solution worse than the best constructive heuristic and improves on it in $25.8\%$ and $43.0\%$ of the
instances. Against the certified bound, $74.7\%$ of its solutions are provably optimal (MULTIFIT: $48.5\%$; LDM: $34.0\%$; \texttt{ig\_ql} started from the same solution: $69.3\%$), and the mean GAP is $0.133\%$ against $0.425\%$ for MULTIFIT.
Tuning $d$ is worth $0.034\%$ (synthetic) and $0.100\%$ (published) over the default $d$, concentrated in $50\leq n\leq75$; for $n\geq150$ the constructive heuristics are already hard to improve.

% Auto-generated by analysis/make_tables_cycle3.py (Entrada 125)
\begin{table}[t]
\centering
\caption{Tuned-$d$ Iterated Greedy started from the best constructive solution (\textsc{IG-Tuned}$^{+}$: random destruction, $d/n$ chosen on separate instances, no learning). Synthetic: 1{,}130 instances; published: 1{,}500 (B 400, C 400, E 400, F 300; $d/n$ per family chosen on 553 instances never run). $\bar d$ in \% (negative: \textsc{IG-Tuned}$^{+}$ better); Wilcoxon, Holm over eight.}
\label{tab:f22_igcv}
{\scriptsize\setlength{\tabcolsep}{3.5pt}
\begin{tabular}{@{}llrrr@{}}
\toprule
Set & Comparison & W/L/T & $\bar d$ (\%) & $p_{\mathrm{Holm}}$ \\
\midrule
Synthetic & E1: vs.\ \textsc{IG-QL}$^{+}$ (default $d$, learned strategies) & 111/30/989 & -0.034 & $1\times10^{-6}$ \\
 & E2: vs.\ \textsc{MultiFit} & 378/0/752 & -0.288 & $4\times10^{-63}$ \\
 & E3: vs.\ \textsc{LDM} & 707/0/423 & -1.061 & $1\times10^{-116}$ \\
 & E4: vs.\ best constructive & 292/0/838 & -0.207 & $4\times10^{-49}$ \\
\midrule
Published & E1: vs.\ \textsc{IG-QL}$^{+}$ (default $d$, learned strategies) & 408/127/965 & -0.100 & $3\times10^{-30}$ \\
 & E2: vs.\ \textsc{MultiFit} & 1098/0/402 & -0.849 & $3\times10^{-180}$ \\
 & E3: vs.\ \textsc{LDM} & 845/0/655 & -1.145 & $4\times10^{-139}$ \\
 & E4: vs.\ best constructive & 645/0/855 & -0.420 & $1\times10^{-106}$ \\
\bottomrule
\end{tabular}}
\end{table}


\paragraph{A local search helps only at an equal number of iterations}
Adding the descent of the hybrid to the recipe gives the same lesson as in Section~\ref{sec:f3f4-ablation}. At an equal number of iterations it helps ($-0.036\%$ synthetic, $-0.072\%$ published; $81.2\%$ of
provably optimal solutions against $74.7\%$), but an iteration costs $2.3$ to $82.9$ times more. When the recipe \emph{without} local search receives the same CPU time ($T_{\mathrm{eq}}$ between
$2{,}298$ and $82{,}855$), it is better: the local search loses by $0.034\%$ ($24$ wins against $115$ losses, $p_{\mathrm{Holm}}=3\times10^{-8}$) and the equal-CPU recipe reaches $89.0\%$ of provably
optimal solutions against $81.2\%$ (mean GAP $0.058\%$ against $0.094\%$; Table~\ref{tab:f26_igcvls}, and Supplementary Section~\ref{S-sec:m-recipe}).

% Auto-generated by analysis/make_tables_cycle5.py (Entrada 129)
\begin{table}[t]
\centering
\caption{Does adding the local search of the hybrid to the practical recipe help? Both variants use random destruction with $d/n$ chosen on separate instances and start from the best of LPT, MULTIFIT and LDM. IG+LS adds the descent (inter-swap, reinsert; two passes) after each reconstruction ($T=1000$). \textsc{IG-Tuned}$^{+}$ at equal CPU runs the recipe without local search for $T_{\mathrm{eq}}$ iterations, $T_{\mathrm{eq}}=1000K$ with $K$ the CPU ratio measured per class ($K$ from $2.3$ to $82.9$; $T_{\mathrm{eq}}$ from $2{,}298$ to $82{,}855$). Synthetic: 1{,}200 instances (1{,}130 with a constructive reference); published: 1{,}500. $\bar d$ in \% of the makespan (negative: IG+LS better); Wilcoxon, Holm over five.}
\label{tab:f26_igcvls}
{\scriptsize\setlength{\tabcolsep}{3.5pt}
\begin{tabular}{@{}lrrrrr@{}}
\toprule
Comparison & $N$ & W/L/T & $\bar d$ (\%) & 90\% interval & $p_{\mathrm{Holm}}$ \\
\midrule
F1: IG+LS vs.\ IG ($T=1000$), synthetic & 1{,}200 & 103/21/1076 & -0.036 & [-0.045, -0.028] & $5\times10^{-10}$ \\
F2: IG+LS vs.\ IG at \emph{equal CPU} ($T_{\mathrm{eq}}$ iterations), synthetic & 1{,}200 & 24/115/1061 & +0.034 & [+0.026, +0.042] & $3\times10^{-8}$ \\
F3: IG+LS vs.\ IG ($T=1000$), published & 1{,}500 & 576/36/888 & -0.072 & [-0.081, -0.063] & $1\times10^{-76}$ \\
F4: IG+LS vs.\ best constructive, synthetic & 1{,}130 & 340/0/790 & -0.245 & [-0.269, -0.220] & $5\times10^{-57}$ \\
F5: IG+LS vs.\ best constructive, published & 1{,}500 & 818/0/682 & -0.492 & [-0.534, -0.452] & $8\times10^{-135}$ \\
\midrule
\multicolumn{6}{l}{Synthetic, certified bound (1{,}130 instances): provably optimal / mean GAP (\%)} \\
\multicolumn{6}{l}{\quad best constructive 55.3\% / 0.342; recipe ($T=1000$) 74.7\% / 0.133; recipe + local search 81.2\% / 0.094; recipe at equal CPU 89.0\% / 0.058} \\
\bottomrule
\end{tabular}}
\end{table}


\paragraph{The budget does the rest}
The strongest variant is the recipe with a larger iteration budget. On six classes of the synthetic benchmark it reaches $86.8\%$, $92.7\%$ and $\mathbf{97.5\%}$ of provably optimal solutions at $T=4{,}000$,
$16{,}000$ and $64{,}000$, against $75.5\%$ at $T=1{,}000$ (mean GAP $0.099\%\to0.009\%$; Table~\ref{tab:s_recipe_budget}); the results at $T=1{,}000$ throughout the paper are comparisons at a fixed, small
budget, not claims about what each method can reach. Exact methods remain the right tool for the instances they solve in about a second (about $85\%$ of families C and F; \citealp{OstojicEtAl2025}); the point
of the recipe is that a learning-free method already sits near the optimum.

% Auto-generated by analysis/make_tables_cycle5.py (Entrada 129)
\begin{table}[t]
\centering
\caption{Quality against budget of the learning-free recipe (Iterated Greedy with random destruction, $d/n$ validated per class, start from the best of LPT, MULTIFIT and LDM, no local search): mean GAP (\%) to the certified bound and share of provably optimal solutions (\%), 100 instances per class. Trajectories share the seed, so the result is non-increasing in $T$ by construction (0 violations). Mean CPU time per run at $T=64{,}000$: $12.6$--$24.3$\,s (parallel run, indicative).}
\label{tab:s_recipe_budget}
{\scriptsize\setlength{\tabcolsep}{3.5pt}
\begin{tabular}{@{}lrrrrrrr@{}}
\toprule
Class & GAP $T=250$ & GAP $T=1{,}000$ & GAP $T=4{,}000$ & GAP $T=16{,}000$ & GAP $T=64{,}000$ & Optimal $T=1{,}000$ & Optimal $T=64{,}000$ \\
\midrule
30x8 & 0.396 & 0.160 & 0.068 & 0.020 & 0.010 & 69.0 & 98.0 \\
50x10 & 0.265 & 0.130 & 0.038 & 0.015 & 0.000 & 67.0 & 100.0 \\
75x15 & 0.207 & 0.138 & 0.084 & 0.044 & 0.008 & 65.0 & 98.0 \\
100x20 & 0.171 & 0.151 & 0.106 & 0.071 & 0.035 & 61.0 & 91.0 \\
150x15 & 0.016 & 0.016 & 0.014 & 0.012 & 0.002 & 92.0 & 99.0 \\
250x25 & 0.002 & 0.002 & 0.002 & 0.002 & 0.002 & 99.0 & 99.0 \\
pooled & 0.176 & 0.099 & 0.052 & 0.027 & 0.009 & 75.5 & 97.5 \\
\bottomrule
\end{tabular}}
\end{table}


\subsection{External validation on published identical-machine instances}
\label{sec:external}

To test whether the findings hold on instances this project did not design, a 13-method roster was run on $6{,}780$ instances of the benchmark released with \citet{OstojicEtAl2025} (a stratified sample of $2{,}000$ from
each of families B, C and E, and all $780$ of F), analysed with the same Wilcoxon/Holm procedure.\footnote{Wide-range instances of family B ($U[1,10^{18}]$ at the extreme) overflow 64-bit load accumulators; the code
switches to arbitrary-precision integers above a safe threshold, with no change on any other instance (verified by independent recomputation and regression tests).}

\paragraph{What replicates, and what does not}
The comparison with the published baselines replicates: \texttt{ig\_ql} against SA, BCO and MULTIFIT is Holm-significant and favourable in a large majority of tests in all four families. The contrast between
mechanisms does not replicate uniformly: the ILS wins nearly all significant comparisons in family B ($157$ of $163$ favour the ILS; wide-range, almost duplicate-free times), the IG nearly all of C ($100$ of $102$), and
E and F are mixed. Neither mechanism dominates; the granularity of the processing times matters more than either fixed choice (Supplementary Section~\ref{S-sec:m-granularity}).

\paragraph{The small policy effect is real and replicates}
A family-level pooled test suggested that \texttt{ig\_ql} has a slightly lower GAP than \texttt{ig\_random} in B, C and E. Because the hypothesis arose after looking at the data, it was tested on $4{,}160$ instances never
run before, with the hypothesis, the test and the decision rule fixed in writing, with hashes, before any of them was run: the rule was met in all three families ($-3.0\%$ and $-3.5\%$ of the GAP in B and E; $-0.85\%$
in C, with an interval that includes zero). A second test on $2{,}690$ further instances (Table~\ref{tab:f16_confirm2}) replicates it in all three families ($0.015$, $0.029$ and $0.027$ percentage points of GAP;
$p_{\mathrm{Holm}}\leq0.003$), and an informative state reproduces it without adding to it. The effect is therefore real, independent of the state and negligible in size: hundredths of a point.

% Auto-generated by analysis/make_tables_cycle2.py (Entrada 122)
\begin{table}[t]
\centering
\caption{Second confirmatory test on 2{,}690 published instances never run before (B 980; C 1{,}050 from 75 classes not sampled earlier; E 660), with the original state (\texttt{ig\_ql}) and with an informative state (\texttt{ig\_ql\_rel}; thresholds calibrated on 553 disjoint instances). $\Delta$GAP: mean difference against \texttt{ig\_random} in percentage points (negative: Q-Learning better); Holm over six tests.}
\label{tab:f16_confirm2}
{\scriptsize\setlength{\tabcolsep}{4pt}
\begin{tabular}{@{}llrrrr@{}}
\toprule
Family & State & $N$ & W/L/T & $\Delta$GAP (pp) [\% of GAP] & $p_{\mathrm{Holm}}$ \\
\midrule
B & original & 980 & 507/348/125 & -0.0147 [-5.4\%] & $2\times10^{-7}$ \\
B & informative & 980 & 493/367/120 & -0.0079 [-2.9\%] & 0.002 \\
\midrule
C & original & 1050 & 268/182/600 & -0.0294 [-1.8\%] & 0.003 \\
C & informative & 1050 & 284/177/589 & -0.0497 [-3.1\%] & $5\times10^{-5}$ \\
\midrule
E & original & 660 & 133/88/439 & -0.0268 [-3.2\%] & 0.002 \\
E & informative & 660 & 130/78/452 & -0.0279 [-3.3\%] & $9\times10^{-4}$ \\
\bottomrule
\end{tabular}}
\end{table}


\paragraph{Gap to exact solutions}
The GAP against the classical bound mixes algorithmic error with bound slack. For families C and F the benchmark provides the solutions of an exact method; \texttt{ig\_ql} is $1.13\%$ above the optimum in C
(optimal in $21.3\%$ of the instances, against $1.68\%$ and $18.3\%$ for MULTIFIT) and $0.17\%$ above in F. The classical bound is loose by $0.72\%$ (C) and $1.01\%$ (F), so a sizeable part of the GAP reported
elsewhere is bound slack, and, conversely, the best heuristic of the study in C is not near-optimal there ($79\%$ of its solutions exceed the optimum).
